\documentclass[aps,prd,twocolumn,showpacs,preprintnumbers,superscriptaddress]{revtex4-2}
\usepackage{amsmath,amssymb}
\usepackage{bm}
\usepackage{booktabs}
\usepackage{array}
\usepackage{tikz}
\usetikzlibrary{calc,arrows.meta}

\newcommand{\SU}{\mathrm{SU}}
\newcommand{\UU}{\mathrm{U}}
\newcommand{\SO}{\mathrm{SO}}
\newcommand{\Nc}{N_c}
\newcommand{\Nf}{N_f}
\newcommand{\Tr}{\mathrm{Tr}}
\newcommand{\adj}{\mathrm{adj}}
\newcommand{\MeV}{\,\mathrm{MeV}}
\newcommand{\GeV}{\,\mathrm{GeV}}
\newcommand{\TeV}{\,\mathrm{TeV}}
\newcommand{\Lqcd}{\Lambda_{\mathrm{QCD}}}
\newcommand{\LF}{\Lambda_F}
\newcommand{\fpi}{f_\pi}

\begin{document}

\title{The sBootstrap:\\
  emergent $\mathcal{N}{=}2$ from QCD confinement, three
  hypermultiplets, and the Standard Model}

\author{A.\ Rivero}
\affiliation{Independent researcher.
  \texttt{rivero@unizar.es}}

\date{\today}

\begin{abstract}
We present the sBootstrap as a complete theory in which all
Standard Model fermions and their scalar partners arise from
a single confining gauge group $\SU(3)_F$ (family symmetry)
with three flavours.
The confined phase produces three $\mathcal{N}{=}2$
hypermultiplets---one per generation---each containing a
charged lepton (mesino), its meson superpartner, a quark
(diquarkino), and its diquark superpartner, all connected
by a single composite supercharge~$Q$.
In the democratic limit (Cartan angle $\delta \to 0$) the
three hypermultiplets are degenerate at the constituent quark
mass $\mu^2 = 314\MeV$; the physical hierarchy is generated
by the vacuum alignment angle $\delta = 2/9$.
SUSY is broken by QCD chiral dynamics: $m_\pi^2 - m_\mu^2 =
f_\pi^2$ for the light sector and $m_\tau = f_\pi m_b/\Lqcd$
for the heavy sector, with the Volkov--Akulov goldstino
identified as the muon.
The Koide formula $K = 2/3$ is a Casimir identity holding iff
$N = 3$ families; the Cabibbo angle equals the Koide angle
($\theta_C = \delta$, algebraically); and all six CKM + PMNS
parameters follow from the $1/N$ expansion:
$|V_{us}| = 2/9$, $|V_{cb}| = 4/99$, $A = 9/11$,
$\theta_{23}^{\mathrm{PMNS}} = \pi/4 + \delta/3$
($0.08\%$ off).
The bion relation $z_b = 3z_s + z_c$ keeps the quark
charge algebra in $\mathbb{Q}[\sqrt{3}]$ and predicts
$m_t = 171.1\GeV$ ($0.8\%$ off).
Fifteen SM parameters follow from three inputs
($N$, $\mu$, $m_c$) at $1.1\%$ mean accuracy.
\end{abstract}

\maketitle

%======================================================================
\section{Introduction}
\label{sec:intro}
%======================================================================

Miyazawa proposed the first hadronic supersymmetry in
1966~\cite{Miyazawa:1966}, pairing mesons $(\pi, K)$ with
baryons $(N, \Lambda)$ in $\SU(6)$ supermultiplets without
Lorentz generators---five years before Gol'fand and Likhtman
formulated the super-Poincar\'e algebra and eight years before
Wess and Zumino linearised SUSY into renormalisable Yang--Mills.
Volkov and Akulov~\cite{VolkovAkulov:1973} showed that a
massless goldstino with decay constant~$f$ acquires mass
$m \sim f^2/M$ from symmetry breaking at scale~$M$---precisely
the structure of $m_\mu \sim f_\pi^2/\Lqcd$ in the sBootstrap.
The idea that mesons and leptons could be superpartners was
abandoned not because of any algebraic obstruction but because
the Standard Model was constructed with fundamental fermions
and composite hadrons as separate entities.
We argue that the separation was premature.
The logic of the Standard Model---that quarks are
``fundamental'' while pions are ``composite''---is an
organisational choice, not a theorem.  Confinement makes
the distinction operational only above the QCD scale; below
it, what one calls fundamental is a matter of which
variables one writes the Lagrangian in.  Seiberg
duality~\cite{Seiberg:1995} makes this explicit: the same
physics admits both an ``electric'' description (in terms of
quarks) and a ``magnetic'' description (in terms of mesons
and baryons), with neither more fundamental than the other.
The sBootstrap takes Seiberg's lesson seriously and writes
the low-energy Lagrangian in the magnetic variables, where
the mesons and diquarks are the elementary fields and the
quarks and leptons are their fermionic partners.

The empirical evidence is striking.  The pion and muon masses
satisfy $m_\pi^2 - m_\mu^2 = f_\pi^2$ to~$2\%$, where $f_\pi =
92\MeV$ is the pion decay constant.  The three charged lepton
masses satisfy Koide's formula~\cite{Koide:1983}
$K = (m_e + m_\mu + m_\tau)/(\sqrt{m_e}+\sqrt{m_\mu}+\sqrt{m_\tau})^2
= 2/3$ to one part in~$10^5$.  The tau mass obeys
$m_\tau = f_\pi\,m_b/\Lqcd$ to~$0.2\%$.  And the meson triple
$(-\pi, D_s, B)$ satisfies its own Koide relation to~$0.1\%$.
These are relations between \emph{different sectors} of the
Standard Model---leptons, mesons, and quarks---that have no
explanation unless the sectors are connected by a symmetry.

The sBootstrap~\cite{Rivero:2024} provides the connection.
It identifies the Standard Model fermions as the fermionic
composites of a confining $\SU(3)_F$ family gauge
theory~\cite{MasieroVeneziano:1985}.
A composite supercharge~$Q$ acts in two sectors:
\begin{align}
  Q_\alpha\,|\text{meson}\rangle &= |\text{lepton}_\alpha\rangle
  &&\text{(colour singlet)}\,,
  \label{eq:Q_meson}\\
  Q_\alpha\,|\text{diquark}\rangle &= |\text{quark}_\alpha\rangle
  &&\text{(colour triplet)}\,.
  \label{eq:Q_diquark}
\end{align}
The mesons and diquarks of QCD are the scalar partners
(sleptons, squarks) of the MSSM---not new particles, but
\emph{the same particles already observed} as hadrons.
The counting was established in Ref.~\cite{Rivero:2024}:
with five light quarks ($u,d,s,c,b$; the top does not
hadronize), the combinatorics of $q\bar q$ and $qq$ pairs
reproduces the scalar content of the MSSM with exactly
three generations.

The composite supercharge carries spin~$1/2$, electric
charge~$0$, colour singlet, baryon number~$0$.
No gauge quantum number obstructs the identification;
``lepton number'' is the fermion number within the
supermultiplet, just as a selectron carries lepton number
in the MSSM.
However, $Q$ cannot be a local QCD operator.
Chadha and Nielsen~\cite{ChadhaNielsen:1977} proved that
soft-pion IR divergences cause the local anticommutator
$\{Q,\bar Q\}_{\mathrm{local}}$ to overshoot
$2\sigma^\mu P_\mu$ by a factor
$\sim N_c^3(m_\pi/\Lqcd)^2 \approx 18$.
The explicit calculation (given in Ref.~\cite{v159}) inserts
the QCD condensate $\langle\bar q q\rangle = -B_0 f_\pi^2$
into the composite operator and confirms the $N_c^3$ scaling.
Three evasions survive: (i)~holographic AdS/QCD with a
warp-factor suppression matching $1/18$~\cite{Karch:2002},
(ii)~Wilson-loop smeared operators over paths
$\sim 5\Lqcd^{-1}$, and (iii)~Volkov--Akulov IR regulation
via the constraint $\Phi_{\mathrm{VA}}^2 = 0$.
Each is testable on the lattice within 3--5 years
(Sec.~\ref{sec:tests}).
In this paper we adopt the Volkov--Akulov realisation as the
working framework: the nonlinear constraint
$\Phi_{\mathrm{VA}}^2 = 0$ removes the scalar component from
the goldstino chiral superfield, eliminating the soft-pion IR
poles that drive the Chadha--Nielsen overshoot.
All three evasions converge on the same physical picture---the
composite supercharge is an effective infrared object, not a
local Noether current.

The sBootstrap differs from Miyazawa's hadronic
supersymmetry~\cite{Miyazawa:1966} in a fundamental way.
Miyazawa pairs mesons with baryons
($\pi \leftrightarrow N$ in the same multiplet),
which requires the supercharge to carry baryon number.
The sBootstrap pairs mesons with leptons
($\pi \leftrightarrow \mu$): the supercharge carries zero
baryon number, and ``lepton number'' is simply the fermion
number within the supermultiplet.
This distinction matters because Miyazawa-type supercharges
fail at the quantitative level (the meson--baryon mass
splittings do not match any known SUSY-breaking pattern),
while the meson--lepton splittings satisfy the Volkov--Akulov
relation $\Delta m^2 = f_\pi^2$ to~$2\%$.
The empirical success of the sBootstrap identification---not
just the algebraic possibility---is what motivates this paper.

In this paper we work with the algebraic and dynamical
consequences of~(\ref{eq:Q_meson})--(\ref{eq:Q_diquark}),
building the complete theory from the ground up.
We describe the particle content (Sec.~\ref{sec:content}),
the three $\mathcal{N}{=}2$ hypermultiplets and their
degeneracy structure (Sec.~\ref{sec:hyper}),
the SUSY-breaking regimes (Sec.~\ref{sec:breaking}),
the Koide formula (Sec.~\ref{sec:koide}),
the quark mass chain (Sec.~\ref{sec:chain}),
the CKM and PMNS mixing (Sec.~\ref{sec:mixing}),
the underlying $\mathcal{N}{=}2$ structure (Sec.~\ref{sec:N2}),
and the M-theory embedding (Sec.~\ref{sec:Mtheory}).
A complete table of~15 predictions at $1.1\%$ mean accuracy
is given in Sec.~\ref{sec:predictions}, followed by
experimental tests in Sec.~\ref{sec:tests}.

%======================================================================
\section{Particle content}
\label{sec:content}
%======================================================================

The central combinatorial fact of the sBootstrap is this:
with $q$ up-type and $r$ down-type light quarks, the number
of charged $q\bar q$ mesons is $qr$ and the number of
antisymmetric scalar diquarks is $\binom{q+r}{2}$.
Requiring that these fill the scalar content of the MSSM
for $n$ generations (with $2n$ sleptons and $3n$ squarks
per colour) and that $r = 2q-1$ (equal numbers of
charge-$+1/3$ mesons from each type), the unique minimal
solution is $q = 2$, $r = 3$, $n = 3$: two up-type quarks
($u, c$), three down-type ($d, s, b$), and exactly three
generations of leptons.
The top quark does not participate because it decays before
confining; its absence is what makes the counting close.
This is the derivation of three
generations~\cite{Rivero:2024}.

\paragraph{Charged mesons: the slepton sector}

Two up-type quarks ($u, c$) paired with three down-type
($d, s, b$) give six charged mesons:
\begin{equation}
  \pi^- = \bar u d,\; K^- = \bar u s,\; B^- = \bar u b,\;
  D^- = \bar c d,\; D_s^- = \bar c s,\; B_c^- = \bar c b.
  \label{eq:6mesons}
\end{equation}
These are identified with the six charged sleptons of the MSSM
(three left, three right)---one scalar per lepton chirality.
The two up-type quarks map to the two chiralities:
$u$-type mesons correspond to $\SU(2)_L$ doublet sleptons
and $c$-type mesons to singlet sleptons.
The quantum numbers of each meson match its proposed
slepton partner once the SU(2)$_L$ assignment is made:

\begin{table}[h]
\centering
\caption{Quantum numbers of the six charged mesons and
  their slepton identifications.  Colour is singlet
  throughout; $I_3$ is the weak isospin projection after
  electroweak embedding.}
\label{tab:QN}
\begin{tabular}{lcccc}
\toprule
Meson & $Q$ & $Y$ & $I_3$ & Slepton \\
\midrule
$\pi^-=\bar ud$ & $-1$ & $-1$ & $-1/2$ & $\tilde e_L$ \\
$K^-=\bar us$ & $-1$ & $-1$ & $-1/2$ & $\tilde\mu_L$ \\
$B^-=\bar ub$ & $-1$ & $-1$ & $-1/2$ & $\tilde\tau_L$ \\
$D^-=\bar cd$ & $-1$ & $-1$ & $0$ & $\tilde e_R$ \\
$D_s^-=\bar cs$ & $-1$ & $-1$ & $0$ & $\tilde\mu_R$ \\
$B_c^-=\bar cb$ & $-1$ & $-1$ & $0$ & $\tilde\tau_R$ \\
\bottomrule
\end{tabular}
\end{table}

The $I_3$ assignment follows from the preon quantum numbers
of Masiero and Veneziano~\cite{MasieroVeneziano:1985}:
the $u$-quark preon carries $I_3 = -1/2$ (doublet) while
the $c$-quark preon is an $\SU(2)_L$ singlet ($I_3 = 0$).
QCD colour is carried as a spectator charge---the
$\SU(3)_F$ confinement does not affect it, which is why
all six mesons are colour singlets and all six sleptons
are colour singlets.

\paragraph{Scalar diquarks: the squark sector}

Scalar (spin-0) diquarks in the colour-$\bar 3$ channel
require flavour antisymmetry:
$\binom{5}{2} = 10$ combinations.
These are the squarks.
The six diquarks with charge $\pm 1/3$ match the
down-type squarks; three with charge~$+1/3$ from $[uc]$,
$[dc]$, $[sc]$ and $[ud]$, $[us]$, $[ds]$ match the
up-type; and one charge-$4/3$ diquark $[uc]$ is the
problematic extra~\cite{Rivero:2024}.

The detailed counting is as follows.
The 10 antisymmetric diquarks
carry charges $+1/3$ (6~combinations: $[ud]$, $[us]$, $[ub]$,
$[ds]$, $[dc]$, $[sc]$), $-1/3$ (3~combinations: $[ds]$,
$[db]$, $[sb]$), and $+4/3$ (1~combination: $[uc]$).
The six charge-$\pm 1/3$ diquarks map to the down-type
squarks ($\tilde d_L, \tilde s_L, \tilde b_L$ and their
right-handed partners); the three remaining charge-$+1/3$
map to up-type squarks.
The single charge-$4/3$ diquark $[uc]$ is the problematic
extra---it has no SM fermion partner.
In the $\SO(32)$ heterotic embedding~\cite{Rivero:2024},
the anomaly-free completion absorbs this state together
with $442$ others from the adjoint decomposition
$\mathbf{496} = \mathbf{54} \oplus \mathbf{442}$
under $\SU(5)_f \subset \SO(32)$.

The 10~diquarks and 6~mesons together give $16$ scalars.
The MSSM has $18$ squarks + sleptons for three generations.
The shortfall of two reflects the absence of the top from
the confined sector: no $[ut]$ or $[ct]$ diquark exists
because the top decays before confining.
The charge-$4/3$ excess is the single diquark $[uc]$,
which has no SM fermion partner at this charge.

The full scalar content is summarised in Table~\ref{tab:scalars}.
The ``good'' sector (charges $\pm 1/3$ and $\pm 1$) matches
the MSSM exactly.
The ``exotic'' charge-$4/3$ sector must either decouple at
high energies (via the $\SO(32)$ projection) or form a
neutral condensate.

\begin{table}[h]
\centering
\caption{Scalar diquarks from five quarks ($u,d,s,c,b$)
  in the colour-$\bar 3$ antisymmetric channel.}
\label{tab:scalars}
\begin{tabular}{lcl}
\toprule
Charge & Count & Flavour content \\
\midrule
$+1/3$ & $6$ & $[ud],[us],[ub],[dc],[sc],[cb]$ \\
$-1/3$ & $3$ & $[ds],[db],[sb]$ \\
$+4/3$ & $1$ & $[uc]$ \\
\midrule
Total & $10$ & \\
\bottomrule
\end{tabular}
\end{table}

\paragraph{Neutral mesons: no lepton partners}

The sBootstrap contains only chiral superfields
($\Phi = (\phi, \psi, F)$) pairing mesons with leptons---no
vector superfields.
The D-terms that normally split charged from neutral
members of a gauge multiplet are absent.
Consequently, neutral mesons
($\pi^0, \eta, K^0, D^0, B^0, B_s$) have
\textbf{no lepton partners}.
Neutrinos arise from a different sector
(Sec.~\ref{sec:neutrinos}).

%======================================================================
\section{Three hypermultiplets}
\label{sec:hyper}
%======================================================================

\paragraph{The $\SU(3)_F$ framework}

The confining gauge group is $\SU(3)_F$ (family symmetry) with
$\Nf = \Nc = 3$ hypermultiplets in the fundamental.
The choice of $\mathcal{N}{=}2$ rather than $\mathcal{N}{=}1$
in the UV is not aesthetic but forced by the Koide formula:
$\mathcal{N}{=}2$ protects the K\"ahler potential from
renormalisation, ensuring that the tree-level mass
relations survive quantum corrections (the wave-function
renormalisation $\delta Z$ is suppressed to
$\sim 10^{-9}$~\cite{v159}).
Without this protection, the Koide formula would receive
$O(\alpha)$ corrections that spoil its $10^{-5}$ precision.
In the UV, the $\mathcal{N}{=}2$ SQCD theory has matter
content $H_i = (Q_i, \tilde Q_i)$ with $i = 1,2,3$.
After s-confinement ($\Nf = \Nc$), the IR degrees of freedom
are the $9$~mesons $M_{ij} = Q_i \tilde Q_j$ (scalars, colour
singlet), the $9$~mesinos $\psi_{M_{ij}}$ (fermions, identified
with leptons), and the baryons
$B = \epsilon_{ijk} Q_i Q_j Q_k$,
$\tilde B = \epsilon_{ijk} \tilde Q_i \tilde Q_j \tilde Q_k$,
subject to the quantum-modified constraint
$\det M - B\tilde B = \Lambda_F^6$.

The 9~mesinos decompose as
$3L_i + 3e^c_i + 3\nu^c_i$ (or equivalently,
$3 \times 3$ under the residual $\SU(3)_L \times \SU(3)_R$
flavour symmetry).
The three diagonal mesinos $\psi_{M_{ii}}$ are the
three charged leptons $(e, \mu, \tau)$.
The six off-diagonal mesinos $\psi_{M_{ij}}$
($i \neq j$) are massless at tree level---they carry the
seeds of CKM and PMNS mixing.

\paragraph{What is in each hypermultiplet}

The three mass eigenvalues of the meson matrix define the
three hypermultiplet levels.
In the Koide parametrisation (Sec.~\ref{sec:koide}),
$\sqrt{m_n} = \mu(1 + \sqrt{2}\cos(\delta + 2\pi n/3))$
with $\delta = 2/9$.

\begin{table}[h]
\centering
\caption{Content of the three hypermultiplets.
  Each pairs a lepton with mesons (colour-singlet sector)
  and quarks with diquarks (colour-triplet sector).
  The supercharge~$Q$ connects each row.}
\label{tab:hyper}
\begin{tabular}{lcccc}
\toprule
Level & Lepton & Mesons & Quarks & Diquarks \\
\midrule
$H_\tau$ & $\tau^-$ & $B^-,\;B_c^-$ & $b,\,t$ &
  $[ub],[cb],[sb]$ \\
$H_\mu$ & $\mu^-$ & $K^-,\;D_s^-$ & $s,\,c$ &
  $[us],[cs],[ds]$ \\
$H_e$ & $e^-$ & $\pi^-,\;D^-$ & $u,\,d$ &
  $[ud]$ \\
\bottomrule
\end{tabular}
\end{table}

The assignment of mesons to generations is \emph{not} by
mass ordering but by the superpotential Yukawa
couplings~$y^\alpha_{ij}$ that determine which meson
F-term feeds which lepton mass (Sec.~\ref{sec:breaking}).

\paragraph{The $\mathcal{N}{=}2$ representation content.}

In the UV $\mathcal{N}{=}2$ SQCD, each hypermultiplet
$H_i = (Q_i, \tilde Q_i)$ contains two $\mathcal{N}{=}1$
chiral multiplets:
\begin{equation}
  Q_i = (\phi_i, \psi_{Q_i})\,,\quad
  \tilde Q_i = (\tilde\phi_i, \psi_{\tilde Q_i})\,,
\end{equation}
where $\phi_i$ are complex scalars (squarks of $\SU(3)_F$)
and $\psi_{Q_i}, \psi_{\tilde Q_i}$ are Weyl fermions.
The $\SU(2)_R$ R-symmetry rotates $(\psi_{Q_i},
\psi_{\tilde Q_i}^\dagger)$ as a doublet.

After $\SU(3)_F$ s-confinement, the composite mesons
$M_{ij} = Q_i \tilde Q_j$ inherit both scalar and
fermion components.
The diagonal mesons $M_{ii}$ form the mass eigenstates
(one per generation), while the off-diagonal $M_{ij}$
($i \neq j$) mediate inter-generational transitions.

In addition to the mesons, the $\mathcal{N}{=}2$ vector
multiplet of $\SU(3)_F$ contains an adjoint scalar~$\Phi$
(8~real components) and adjoint fermions (gauginos).
The UV-to-IR chain proceeds in three steps:
(i)~$\mathcal{N}{=}2$ SQCD with the superpotential
$W = \sqrt{2}\,g\,\tilde Q_i\Phi Q^i$ at high energies;
(ii)~$\mathcal{N}{=}2 \to \mathcal{N}{=}1$ breaking by an
adjoint mass $\mu_\Phi\,\Tr(\Phi^2)/2$, which after
integrating out~$\Phi$ gives
\begin{equation}
  W_{\mathrm{eff}} = -\frac{g^2}{\mu_\Phi}\Bigl[\Tr(M^2)
  - \frac{1}{3}(\Tr M)^2\Bigr] + m_i M^i_i\,;
  \label{eq:Weff}
\end{equation}
and (iii)~$\SU(3)_F$ s-confinement replacing the elementary
fields with mesons and baryons subject to
$\det M - B\tilde B = \Lambda_F^6$.
The adjoint mass
$\mu_\Phi = 11\,m_\tau/2 \approx 9.8\GeV$
(Sec.~\ref{sec:N2}) is well above the confinement scale
$\Lambda_F \sim 100\MeV$, so steps~(ii) and~(iii) are
cleanly separated.
The gaugino condensate contributes to the SUSY breaking
and generates the angular potential for~$\delta$.

Each hypermultiplet level therefore contains, after
confinement and SUSY breaking:
(i)~one charged lepton (the mesino fermion),
(ii)~two charged mesons (the $\SU(2)_R$ doublet scalars,
  one per up-type quark),
(iii)~one or more quarks (the diquarkino fermions),
(iv)~the corresponding scalar diquarks.
Table~\ref{tab:hyper} summarises the assignment.

\paragraph{The democratic limit.}

At $\delta = 0$, all three levels are degenerate:
\begin{equation}
  m_1 = m_2 = m_3 = \mu^2 = 314\MeV\,.
  \label{eq:democratic}
\end{equation}
This is the \textbf{constituent quark mass}---the scale of
QCD confinement, not a free parameter.
The degeneracy reflects the unbroken $\SU(3)_F$ flavour
symmetry and is \emph{not} an enhancement to
$\mathcal{N} > 2$: the $\SU(2)_R$ of $\mathcal{N}{=}2$
acts \emph{within} each hypermultiplet (relating $L$ to
$e^c$), while $\SU(3)_F$ acts \emph{between} them.
No supercharge can relate different generations because
they carry different $\SU(3)_F$ quantum numbers.

There is no physical $2+1$ degeneracy.  The Koide
parametrisation gives three distinct masses for any
$\delta \neq 0$.  However, the physical spectrum is
\emph{approximately} $2+1$ since $m_e \ll m_\mu \ll m_\tau$;
the exact $2+1$ point is the seed $\delta = \pi/12$ where
one mass vanishes exactly.

\paragraph{$\SU(2)_R$ and the L/R split}

The $\SU(2)_R$ of $\mathcal{N}{=}2$ rotates
$(Q_i, \tilde Q_i^\dagger)$ within each hypermultiplet.
After confinement:
\begin{equation}
  \underbrace{\psi_{Q_i}}_{\displaystyle L_i}
  \;\xleftrightarrow{\SU(2)_R}\;
  \underbrace{\psi_{\tilde Q_i}^\dagger}_{\displaystyle e_i^c}
  \label{eq:SU2R}
\end{equation}
The left--right split of the Standard Model---that $L$
transforms as an $\SU(2)_L$ doublet while $e^c$ is a
singlet---is the spectral remnant of $\SU(2)_R$ breaking
when $\mathcal{N}{=}2 \to \mathcal{N}{=}1$.

The two mesons per generation (one $u$-type, one $c$-type)
are the two scalar components of the $\SU(2)_R$ doublet:
\begin{equation}
  (\pi^-, D^-) = \SU(2)_R\;\text{doublet of }H_e\,,
  \label{eq:doublet}
\end{equation}
and similarly $(K^-, D_s^-)$ for $H_\mu$ and
$(B^-, B_c^-)$ for $H_\tau$.

%======================================================================
\section{SUSY breaking: three regimes}
\label{sec:breaking}
%======================================================================

The key physical insight of the sBootstrap is that
\textbf{chiral symmetry breaking IS the SUSY breaking}.
The two are not independent phenomena occurring at different
scales; they are the same phase transition, with order
parameter $\langle\bar q q\rangle \neq 0$ and breaking
scale $f = f_\pi = 92\MeV$.
This identification resolves the ``SUSY flavour problem''
that plagues the MSSM: there is no separate hidden sector,
no messenger mechanism, no mediation scale to tune.
The goldstino is the muon, its decay constant is $f_\pi$,
and the SUSY-breaking scale is the QCD scale.

The Volkov--Akulov~\cite{VolkovAkulov:1973} nonlinear
realisation makes this precise.
In standard (linear) SUSY breaking, soft masses satisfy
$\Delta m^2 \propto |F|^2/M_{\mathrm{Pl}}^2$---gravity-mediated
and Planck-suppressed.
In the sBootstrap, the splitting is \emph{first-order}
in~$\langle F\rangle$: $\Delta m^2 \propto f_\pi^2$, with
no Planck suppression.
This is the signature of nonlinear (strong) SUSY breaking,
where the goldstino is a QCD composite and the breaking
is mediated by strong dynamics, not gravity.

The splitting depends on the quark mass relative
to~$\Lqcd$, giving rise to three distinct regimes.

\paragraph{Regime~I: light quarks ($m_q \ll \Lqcd$).}

For $u, d, s$ quarks, the meson is a pseudo-Goldstone boson
whose mass is set by the GMOR relation
$m_\pi^2 = (m_u + m_d)\,B_0$ with
$B_0 = -\langle\bar q q\rangle/f_\pi^2$.
The meson--lepton splitting is first-order in the
Volkov--Akulov breaking VEV:
\begin{equation}
  m_\pi^2 - m_\mu^2 = f_\pi^2\,.
  \label{eq:regime1}
\end{equation}
Numerically:
$m_{\pi^\pm}^2 - m_\mu^2 = 19480 - 11164 = 8316\MeV^2$
versus $f_\pi^2 = 8477\MeV^2$; ratio $= 0.981$.
The $2\%$ discrepancy lies within the ChPT systematic budget:
Gasser--Leutwyler NLO corrections contribute $+0.2\%$,
the electromagnetic splitting
$m_{\pi^\pm}^2 - m_{\pi^0}^2 = 1261\MeV^2$ adds $6.5\%$,
and lattice $f_\pi$ uncertainty adds~$\pm 1\%$.
The quadrature sum~$\approx 7\%$ encompasses the observed
ratio at~$0.3\sigma$ from
unity---\textbf{no tension with Standard Model QCD}.

The muon mass arises via the PCSC
(partially conserved supercurrent) relation,
by analogy with the PCAC relation
$\partial_\mu A^\mu = f_\pi m_\pi^2 \pi$:
\begin{equation}
  \partial_\mu S^\mu_\alpha(x) = m_\mu\,\bar\mu_\alpha(x)\,,
  \label{eq:PCSC}
\end{equation}
identifying the muon as a quasi-goldstino with decay
constant $f = f_\pi = 92\MeV$.
The dimensional estimate
\begin{equation}
  m_\mu = |y^\mu_{ud}|\,\frac{f_\pi^2}{\Lqcd}
  \approx 40\MeV \times \mathcal{O}(1)\;\text{chiral-loop factors}
  \label{eq:muon}
\end{equation}
reproduces the observed~$105.7\MeV$ with $|y^\mu_{ud}| \approx
5.8 \times 10^{-3}\GeV^{-1}$, consistent with a
nonrenormalisable effective theory below~$\Lqcd$.
The \textbf{linear} scaling $m_\mu \propto f_\pi^2\Lqcd$
(not $f_\pi^2/M_{\mathrm{Pl}}$) is the hallmark of
Volkov--Akulov breaking: the goldstino is a QCD composite,
not a gravity-mediated soft mass.

\paragraph{Regime~II: medium quarks ($m_q \sim \Lqcd$).}

For $c, b$ quarks, explicit quark masses compete with chiral
dynamics.  The meson F-term becomes
$\langle F_{ij}\rangle \sim f_\pi\,m_{q_i q_j}$
for heavy quark pairs (the GMOR relation is replaced by
the heavy-quark expansion), and the splitting interpolates
from $\Delta m^2 \sim f_\pi^2$ to $\Delta m^2 \sim f_\pi m_q$:
\begin{equation}
  m_\tau = \frac{f_\pi\,m_b}{\Lqcd}\,,\qquad
  \Lambda = 216.6\MeV \approx \Lambda_{\overline{MS}}^{(5)}\,.
  \label{eq:regime2}
\end{equation}
Numerically: $92 \times 4183 / 216.6 = 1777\MeV$.
Observed: $1776.86\MeV$.  Accuracy: $0.2\%$.

This second sBootstrap relation connects the tau mass to the
$b$-quark mass through the pion decay constant---a
startling link between the heaviest charged lepton and a
quark that seems unrelated in the Standard Model.
The physical origin is the superpotential coupling
$y^\tau_{ub}\,\Phi_{ub}\,L_\tau$
(Eq.~\ref{eq:superpotential}) with the
$B$-meson F-term $\langle F_{ub}\rangle \sim f_\pi\,m_b$.
The $\Lqcd$ in the denominator reflects the Yukawa coupling
$|y^\tau| \sim 1/\Lqcd$ natural for an effective theory
at the QCD scale.

The pion--muon pairing ($m_\pi^2 - m_\mu^2 = f_\pi^2$) and
the tau--bottom pairing ($m_\tau = f_\pi m_b/\Lqcd$) together
imply \textbf{CKM-like mixing in the meson--lepton sector}:
the $\pi$ pairs with the $\mu$ (not the $e$) and the $B$
pairs with the $\tau$ (not the $\mu$), requiring a non-trivial
Yukawa matrix~$y^\alpha_{ij}$ with off-diagonal structure.

\paragraph{Regime~III: top quark ($m_t \gg \Lqcd$).}

The top decays before hadronizing: its lifetime
$\tau_t \sim 5 \times 10^{-25}\,\text{s}$ is shorter than
the QCD confinement time
$\Lambda_{\mathrm{QCD}}^{-1} \sim 10^{-24}\,\text{s}$.
It forms no meson and therefore has no meson superpartner
in QCD.
The top mass enters only through the algebraic chain
(Sec.~\ref{sec:chain}), not through a direct SUSY pairing.

This is why the sBootstrap uses five quarks ($u,d,s,c,b$),
not six: the top's absence from the confined spectrum is
what makes the combinatorial counting work for exactly three
generations.
With six quarks, the meson count would be $3 \times 3 = 9$
(three up-type, three down-type), giving too many sleptons
for three generations.
The fact that the top does not hadronize is not a blemish
but a \emph{prediction}: it is the mechanism that fixes
the number of generations to three.

\paragraph{The three-regime interpolation.}

The SUSY splitting interpolates smoothly between the regimes:
\begin{equation}
  \Delta m^2 \approx f_\pi^2 + f_\pi\,m_q + \mathcal{O}(m_q^2/\Lqcd^2)\,,
\end{equation}
where the first term dominates for light quarks and the
second for heavy.
The pion--muon relation $m_\pi^2 - m_\mu^2 = f_\pi^2$
(Regime~I) and the tau--bottom relation
$m_\tau = f_\pi m_b/\Lqcd$ (Regime~II) are the leading
terms in this expansion.
The meson Koide triple $K(-\pi, D_s, B) = 0.667$ ($0.1\%$
from $2/3$) is a consistency check: the three mesons that
pair with the three leptons satisfy their own Koide-like
relation, with the negative sign for the pion reflecting
the reversed mass ordering ($m_\pi < m_\mu$ while
$m_{D_s} > m_\mu$ and $m_B > m_\tau$).

\paragraph{The meson--lepton Yukawa matrix}

The superpotential coupling mesons to leptons is
\begin{equation}
  W = \sum_{\alpha=1}^3 \sum_{i,j} y^\alpha_{ij}\,\Phi_{ij}\,L_\alpha
  + \frac{\lambda_{ijk\ell}}{2\Lqcd}\,\Phi_{ij}\Phi_{k\ell}
  + \text{h.c.}
  \label{eq:superpotential}
\end{equation}
where $\Phi_{ij} = (M_{ij}, \psi_{ij}, F_{ij})$ is the
meson chiral superfield and $L_\alpha$ the lepton superfield.
The lepton mass matrix is
\begin{equation}
  m_{\ell,\alpha\beta} = \sum_{ij} y^\alpha_{ij}\,
  \langle F_{ij}\rangle\,(y^\dagger)_{ij\beta}\,.
  \label{eq:lepton_mass}
\end{equation}
The Koide formula constrains the eigenvalues; the Cartan
angle~$\delta$ determines the eigenvectors and hence the
mixing.

The minimal ansatz reproducing
$m_\mu \approx f_\pi$ and $m_\tau \approx f_\pi m_b/\Lqcd$:
\begin{equation}
  y^e \sim 0,\quad
  y^\mu_{ud} \sim g,\quad
  y^\tau_{ub} \sim g\sqrt{m_b/\Lqcd}\,,
  \label{eq:yukawa_ansatz}
\end{equation}
with $g \sim \mathcal{O}(1)$.  The electron mass vanishes
at this order and requires higher-loop or electroweak
corrections.

\paragraph{K\"ahler potential and F-term dynamics.}

The chiral superfield Lagrangian in the confined phase is
$\mathcal{L} = \int d^4\theta\,K(\Phi,\bar\Phi)
+ [\int d^2\theta\,W(\Phi) + \text{h.c.}]$ with K\"ahler
potential
\begin{equation}
  K = \bar\Phi_{ij}\Phi_{ij} + \bar L_\alpha L_\alpha
  + \frac{c_{ijk\ell}}{\Lqcd^2}\,
  (\bar\Phi_{ij}\Phi_{ij})(\bar\Phi_{k\ell}\Phi_{k\ell})
  + \ldots
  \label{eq:kahler}
\end{equation}
At the chiral-breaking minimum, $\langle\Phi_{ij}\rangle = 0$
but the auxiliary field acquires a VEV from the quark
condensate:
$\langle F_\pi\rangle \approx -f_\pi^2\Lqcd$.
The linear scaling $m_\mu \propto f_\pi^2\Lqcd$ (rather than
$f_\pi^2/M_{\mathrm{Pl}}$) is the hallmark of
Volkov--Akulov~\cite{VolkovAkulov:1973} nonlinear breaking,
where the goldstino decay constant is $f = f_\pi$ and the
breaking is mediated by strong dynamics, not gravity.

The constraint $\Phi_{\mathrm{VA}}^2 = 0$ of the
Volkov--Akulov realisation forces the scalar component
to vanish, leaving only the goldstino $\chi$ and the
auxiliary field~$F_{\mathrm{VA}}$.
This eliminates the Chadha--Nielsen soft-pion poles:
the $\chi\chi$ vanishing enforced by the constraint removes
the IR divergence that blocks local composite supercharges.
Gauge invariance requires Wilson lines connecting the quark
bilinears in the meson field to the lepton goldstino.

\paragraph{Neutrinos}
\label{sec:neutrinos}

Neutral mesons have no lepton partners
(Sec.~\ref{sec:content}).  Neutrinos arise from the
\emph{baryon sector}: $B$ and $\tilde B$ of the
$\SU(3)_F$ s-confinement play the role of right-handed
neutrinos.
The quantum-modified constraint
$\det M - B\tilde B = \Lambda_F^6$ links the baryon mass
scale to the meson condensate.
The seesaw mechanism $m_\nu = m_D^T M_R^{-1} m_D$
requires a right-handed mass~$M_R$ much larger than~$\Lambda_F$;
a natural candidate is the inverse Koide crossing scale
$\mu_{\mathrm{cross}} \sim 280\TeV$ where
$K^{-1}(d,s,b) = 2/3$ exactly under 5-loop RG running.
With $M_R \sim N\,\mu_{\mathrm{cross}} \sim 840\TeV$
and $m_D \sim v = 246\GeV$:
$m_\nu \sim v^2/M_R \sim 72\,\mathrm{meV}$,
in the right ballpark for the heaviest neutrino mass.

Assuming a neutrino Koide angle
$\delta_\nu = \delta + \pi/12$
(the seed angle, $= 1/8$ of the $\mathbb{Z}_3$ period
on the Cartan circle), the mass-squared ratio
$\Delta m^2_{32}/\Delta m^2_{21}$ is reproduced and
normal ordering is predicted:
$m_1 \approx 0.35$, $m_2 \approx 8.6$,
$m_3 \approx 51\,\mathrm{meV}$,
$\sum m_\nu \approx 60\,\mathrm{meV}$
(below the Planck bound $120\,\mathrm{meV}$).
The neutrino Koide angle is currently fitted, not derived;
its origin from the seesaw mechanism is an open question.

%======================================================================
\section{The Koide formula from $\SU(3)_F$}
\label{sec:koide}
%======================================================================

The Koide formula~\cite{Koide:1983}
\begin{equation}
  K = \frac{m_e + m_\mu + m_\tau}{(\sqrt{m_e}
  + \sqrt{m_\mu} + \sqrt{m_\tau})^2} = \frac{2}{3}
  \label{eq:koide}
\end{equation}
has held to one part in~$10^5$ since 1981.
In the Brannen parametrisation~\cite{Brannen:2006},
the family charges are
$z_n = 1 + \sqrt{2}\cos(\delta + 2\pi n/3)$
with $\sqrt{m_n} = \mu\,z_n$ and~$\delta$ the Cartan angle.
Two trigonometric identities then give $K = 2/3$
\emph{identically}:
\begin{equation}
  \sum_{n=0}^2 z_n = 3\,,\qquad
  \sum_{n=0}^2 z_n^2 = 6\,,
  \label{eq:trig_identities}
\end{equation}
because $\sum\cos(\delta + 2\pi n/3) = 0$ and
$\sum\cos^2(\delta + 2\pi n/3) = 3/2$ for any~$\delta$.
These are not approximations; they are trigonometric
identities following from the roots-of-unity
structure~$e^{2\pi i n/3}$.
It follows that
$K = \sum z_n^2 / (\sum z_n)^2 = 6/9 = 2/3$,
valid for all~$\delta$ where every $z_n > 0$
(i.e.\ $\delta \in [0, \pi/6)$, which includes
$\delta = 2/9 = 0.222$).

The Koide formula is therefore not a constraint on~$\delta$
but an identity of the parametrisation.  This is not a
weakness---it is the explanation.
The physical content is that the lepton masses \emph{fit}
this parametrisation at all: the mass matrix has the
democratic-plus-Cartan structure of a confined $\SU(3)$
theory with three degenerate flavours split by a single
angular parameter.
In QCD language, the three lepton masses are the three
eigenvalues of a $3 \times 3$ matrix that is ``almost
democratic'' (all entries nearly equal) with a small
perturbation controlled by~$\delta$.
The Koide formula is the statement that the unperturbed
matrix has rank~1 (the democratic matrix
$J_{ij} = 1$), and the perturbation respects the
$\mathbb{Z}_3$ cyclic symmetry of the Cartan subalgebra.

\paragraph{The Casimir interpretation.}

The value $K = 2/3$ has a representation-theoretic meaning.
In the spin-$j$ representation of $\SU(2)$ acting on
$\mathbb{C}^N$ with $j = (N-1)/2$, the expectation value of
$J_3^2$ in the state $|z\rangle = (\sqrt{m_1}, \sqrt{m_2},
\sqrt{m_3})$ (normalised to $\langle z|z\rangle = \sum m_i$)
equals
\begin{equation}
  \langle J_3^2\rangle = \frac{j(j+1)}{N} = K\,.
  \label{eq:casimir}
\end{equation}
For $K = 2/3$ this requires $N(N^2-1) = 24$, which has the
unique solution $N = 3$ among positive integers (verified:
$N=2$ gives $6 \neq 24$; $N=4$ gives $60 \neq 24$; $N=5$
gives $120 \neq 24$).
Physically, the Casimir condition says that the mass vector
$|\sqrt{m}\rangle$ has equal weight in all angular momentum
directions---an isotropy condition that singles out three
families.

This Casimir-Koide identity is one of seven criteria
selecting $N = 3$ families:
(i)~Diophantine ($N(N^2-1)|24$),
(ii)~discriminant ($\mathrm{disc} = 21$),
(iii)~binomial,
(iv)~CP-phase,
(v)~Casimir-Koide ($\langle J_3^2\rangle = K$),
(vi)~PMNS uniqueness ($\chi^2 = 0.5$ only at $N = 3$;
  Sec.~\ref{sec:mixing}), and
(vii)~traceless PMNS ($-1 + 1/N + K = 0$ iff $N = 1$ or~$3$;
  Sec.~\ref{sec:mixing}).
Each is algebraically independent; together they provide
overwhelming evidence that $N = 3$ is the unique number of
families consistent with the $\SU(N)_F$ structure.

\paragraph{Foot's geometric picture.}

Geometrically~\cite{Foot:1994}, the vector
$(\sqrt{m_e}, \sqrt{m_\mu}, \sqrt{m_\tau})$ makes
angle~$\theta$ with the democratic direction $(1,1,1)$:
$\cos\theta = (\sum\sqrt{m_i})/(\sqrt{3}\sqrt{\sum m_i})$.
The Koide formula is equivalent to $\theta = \pi/4$.
The mass vector sits on a cone at $45^\circ$ to the
democratic axis.  The position on the cone is parameterised
by the Cartan angle
\begin{equation}
  \delta = 2/9 = 0.22222\;\text{rad}\,,
  \label{eq:delta}
\end{equation}
experimentally confirmed to better than $10^{-5}$.

\paragraph{The angular potential.}

In the $\SU(3)_F$ framework, $\delta$ parameterises the
vacuum alignment on the Cartan torus.  The $\mathbb{Z}_6$
discrete symmetry of the SU(3) vacuum (from the anomaly
coefficient $2\Nf = 6$) generates a periodic potential.
The first two Fourier harmonics give
\begin{equation}
  V(\delta) = -\cos(3\delta) + \frac{1}{\pi}\cos(6\delta)\,,
  \label{eq:potential}
\end{equation}
The minimum satisfies $\cos(3\delta) = 1/(4\alpha)$ with
$\alpha = 1/\pi$; at this minimum
$\delta_{\min} = 0.2225$, which is $0.12\%$ from~$2/9$.

The coefficient $\alpha = 1/\pi$ has a concrete origin.
The first harmonic $-\cos(3\delta)$ is generated by the
leading instanton of $\SU(3)_F$, which respects the
$\mathbb{Z}_3$ subgroup of the $\mathbb{Z}_6$ anomalous
$\UU(1)_R$.
The second harmonic $+(1/\pi)\cos(6\delta)$ comes from
the one-loop fluctuation determinant around this instanton:
the factor $1/\pi$ is the ratio of the two-instanton to
one-instanton amplitudes, controlled by the $\beta$-function
coefficient $b_0 = 3\Nc - \Nf = 6$.
A full derivation requires computing the determinant of
small fluctuations around the $\SU(3)_F$ instanton at the
confinement scale---a calculation we defer but that is in
principle accessible to lattice methods.

If the $1/\pi$ coefficient is confirmed, $\delta$ becomes
a dynamical prediction: every mass and mixing parameter
follows from $N = 3$ alone, with no free parameters in the
flavour sector.

%======================================================================
\section{The quark mass chain and bion relations}
\label{sec:chain}
%======================================================================

\paragraph{The charge-squared rule}

In the eightfold way, hadron masses are determined by
$\SU(3)$ flavour quantum numbers---the Gell-Mann--Okubo
formula $m = a + bY + c[I(I+1) - Y^2/4]$.
In the sBootstrap, fermion masses are determined by
$\SU(3)_F$ family charges:
$m_f = \mu\,z_f^2$,
where $z_f = 1 + \sqrt{2}\cos(\delta + 2\pi n/3)$ is the
family charge and $\sqrt{m_f} = \mu\,z_f$.
The analogy is precise: just as the Gell-Mann--Okubo formula
follows from the breaking of $\SU(3)_{\mathrm{flavour}}$ by
the strange quark mass, the Koide formula follows from
the breaking of $\SU(3)_F$ by the Cartan angle~$\delta$.
The physical meaning is that $\sqrt{m}$ is the primary
variable---the charge---as supported by the Karasik baryon
current construction~\cite{Karasik:2020}, where the chiral
field $U$ is naturally written as $U = \xi^2$.

\paragraph{The exact seed}

At $\delta_{\mathrm{seed}} = \pi/12$ (one mass exactly
zero), the family charges are algebraic in~$\sqrt{3}$:
\begin{equation}
  z_s = \frac{3-\sqrt{3}}{2}\,,\quad
  z_c = \frac{3+\sqrt{3}}{2}\,,\quad
  \frac{z_c}{z_s} = 2 + \sqrt{3}\,.
  \label{eq:seed_charges}
\end{equation}

\paragraph{Bion relations}

The bion relation~\cite{Unsal:2008}
\begin{equation}
  z_b = 3\,z_s + z_c = 6 - \sqrt{3}
  \label{eq:bion1}
\end{equation}
keeps the system in $\mathbb{Q}[\sqrt{3}]$.
Experimentally, using the standard PDG convention
($m_s$ at $2\GeV$, $m_c$ at $m_c$, $m_b$ at $m_b$---each at
its own natural scale):
$3z_s + z_c = 64.63$ vs $z_b = 64.68$ ($0.07\%$ off).
\emph{Caveat:} at a consistent $\overline{\mathrm{MS}}$ scale
(all at $2\GeV$), the bion is $14\%$ off because $m_c$ and
$m_b$ run significantly between their natural scales and
$2\GeV$.
The $0.07\%$ accuracy is a statement about pole-like masses,
not running masses at a common~$\mu$.

A second relation
\begin{equation}
  z_t = 2\,z_b + 8\,z_c = 24 + 2\sqrt{3}
  \label{eq:bion2}
\end{equation}
predicts $m_t = 171.1\GeV$ ($0.8\%$ from PDG).
The alternative Koide prediction ($K(c,b,t) = 2/3$) gives
$167.2\GeV$ ($3.1\%$ off) and forces the number field into
the degree-4 extension
$\mathbb{Q}[\sqrt{3}, \sqrt{35 - 18\sqrt{3}}]$.
The bion is algebraically simpler and empirically more
accurate.

The algebraic structure has physical significance.
The charges $z_s, z_c, z_b$ in the exact chain
are elements of $\mathbb{Q}[\sqrt{3}]$---a degree-2
extension of the rationals, reflecting the $\SU(3)_F$
structure ($\sqrt{3}$ appears in the Cartan subalgebra
of~$\SU(3)$).
The bion relation preserves this field: $z_b = 6-\sqrt{3}
\in \mathbb{Q}[\sqrt{3}]$.
The Koide prediction for~$z_t$, by contrast, introduces
$\sqrt{35-18\sqrt{3}}$---a degree-4 extension that cannot
be reduced.
The fact that the bion ($\mathbb{Q}[\sqrt{3}]$, degree~2)
outperforms Koide ($\mathbb{Q}[\sqrt{3},\sqrt{35-18\sqrt{3}}]$,
degree~4) suggests that the bion relation is the more
simpler algebraic structure, with the Koide formula
being an approximation valid only when the number field
does not need to extend.

\paragraph{The down-quark system}

With $m_c$ as input, the down-type spectrum is fully determined
by three structural conditions:
$m_d/m_s = \delta^2 = 4/81$ (from $\theta_C = \delta$ combined
with the Gatto--Sartori--Tonin relation~\cite{GattoSartoriTonin:1968}),
the inverse Koide condition $K^{-1}(d,s,b) = 2/3$ (exact at
$280\TeV$ under 5-loop SM RG running), and the bion
relation~(\ref{eq:bion1}).
These three equations in three unknowns have a unique solution.
With $m_c = 1270\MeV$:
\begin{center}
\begin{tabular}{lrrl}
\toprule
& Pred. & PDG & Off \\
\midrule
$m_d$ & $4.5\MeV$ & $4.67$ & $3.6\%$ \\
$m_s$ & $91.6\MeV$ & $93.4$ & $1.9\%$ \\
$m_b$ & $4141\MeV$ & $4183$ & $1.0\%$ \\
$m_t$ & $171.1\GeV$ & $172.6\GeV$ & $0.8\%$ \\
\bottomrule
\end{tabular}
\end{center}

%======================================================================
\section{CKM and PMNS from the Koide angle}
\label{sec:mixing}
%======================================================================

\paragraph{Cabibbo $=$ Cartan.}

The most surprising result of this paper is that the
Cabibbo angle---the oldest and most precisely measured
quark mixing parameter---is algebraically identical to the
Koide angle that governs the lepton mass hierarchy.
This is not a numerical coincidence at the $1\%$ level;
it is an algebraic identity valid for \emph{any} set of
masses satisfying the Koide parametrisation.

Koide's 1981 Cabibbo prediction~\cite{Koide:1983}
\begin{equation}
  \tan\theta_C = \frac{\sqrt{3}(\sqrt{m_\mu} - \sqrt{m_e})}
  {2\sqrt{m_\tau} - \sqrt{m_\mu} - \sqrt{m_e}}
  \label{eq:cabibbo}
\end{equation}
is \emph{algebraically identical} to $\tan\delta$.
To see this, write $\sqrt{m_n} = \mu z_n$ with
$z_n = 1 + \sqrt{2}\cos(\delta + 2\pi n/3)$.
The numerator of~(\ref{eq:cabibbo}) is
$\sqrt{3}(z_2 - z_1) = \sqrt{3}\cdot\sqrt{2}
[\cos(\delta+4\pi/3) - \cos(\delta+2\pi/3)]
= \sqrt{3}\cdot\sqrt{2}\cdot 2\sin\delta\cdot\sin(2\pi/3)
= 2\sqrt{2}\cdot\frac{3}{2}\sin\delta$,
and the denominator is
$2z_0 - z_2 - z_1 = 2\sqrt{2}[\cos\delta -
\frac{1}{2}\cos(\delta+4\pi/3) - \frac{1}{2}\cos(\delta+2\pi/3)]
= 2\sqrt{2}\cdot\frac{3}{2}\cos\delta$.
The ratio is $\tan\delta$, identically for any~$\delta$.
The Cabibbo angle \emph{is} the Koide angle.

\paragraph{The CKM from the democratic resolvent.}

The CKM mixing angle between generations $i$ and $j$ is
determined by the propagation amplitude through the
``generation space'' defined by the democratic matrix
$J_{ij} = 1$ (all entries unity).
The matrix~$J$ has eigenvalue~$3$ (eigenvector $(1,1,1)/\sqrt{3}$,
the democratic direction) and eigenvalue~$0$ (doubly degenerate,
the breaking plane).
The resolvent of the zero-eigenvalue subspace at
$z = 1 + \delta$ is $G(z) = 1/z = 1/(1+\delta)$.
This gives the Wolfenstein $A$ parameter:
\begin{equation}
  A = \frac{1}{1+\delta} = \frac{N^2}{N^2+2}
  = \frac{9}{11}\,,
  \label{eq:A}
\end{equation}
where we used $\delta = 2/N^2$.

The physical interpretation: each adjacent-generation transition
costs a factor~$\delta$ (the ``hopping amplitude''),
and the full CKM element resums the geometric series of
back-and-forth transitions:
$|V_{cb}| = \delta^2 \sum_{n=0}^\infty (-\delta)^n
= \delta^2/(1+\delta)$.
Explicitly:
\begin{equation}
  |V_{us}| = \frac{2}{9}\,,\quad
  |V_{cb}| = \frac{4}{99}\,,\quad
  A = \frac{9}{11}\,.
  \label{eq:CKM}
\end{equation}
Agreement with PDG: $1.2\%$, $4.3\%$, $1.0\%$ respectively.
The $|V_{cb}|$ prediction $4/99 = 0.0404$ lies between the
exclusive determination $0.0395 \pm 0.0008$
($1.1\sigma$ away) and the inclusive $0.0422 \pm 0.0008$
($2.3\sigma$).
Resolution of the long-standing exclusive--inclusive
tension will sharpen this test.

\paragraph{PMNS from the $1/N$ expansion.}

The PMNS mixing angles are deviations from maximal mixing
($\pi/4$), controlled by~$\delta$:
\begin{align}
  \theta_{12} &= \frac{\pi}{4} - \delta + \frac{\delta^2}{N}
  = \frac{\pi}{4} - \frac{2}{9} + \frac{4}{243}\,,
  \label{eq:t12}\\
  \theta_{23} &= \frac{\pi}{4} + \frac{\delta}{N}
  = \frac{\pi}{4} + \frac{2}{27}\,,
  \label{eq:t23}\\
  \theta_{13} &= \frac{2j(j+1)}{N^3} = \frac{4}{27}\;\text{rad}
  = 8.49^\circ\,.
  \label{eq:t13}
\end{align}
The solar angle~(\ref{eq:t12}) satisfies quark--lepton
complementarity~\cite{Raidal:2004} at leading order
($\theta_{12} + \theta_C = \pi/4$) with an NLO correction
$\delta^2/N$.
The atmospheric angle~(\ref{eq:t23}) deviates from maximal
by $\delta/N = 2/27$, a correction of $O(1/N^3)$.
The reactor angle~(\ref{eq:t13}) equals
$K \cdot \delta = (2/3)(2/9) = 4/27$, where $K = 2/3$ is
the Koide number itself.

Agreement with NuFIT~5.2:
$\sin^2\theta_{12} = 0.300$ ($1.3\%$ from $0.304$),
$\sin^2\theta_{23} = 0.574$ ($0.1\%$ from $0.573$),
$\sin^2\theta_{13} = 0.0218$ ($1.9\%$ from $0.0222$).

The correction coefficients $(-1, +1/3, +2/3)$ multiplying
$\delta$ in each angle sum to zero (traceless), reflecting
the $\SU(3)_F$ structure.
The traceless condition $-1 + 1/N + K = 0$ is itself a
constraint on~$N$.
With $K = j(j+1)/N = (N^2-1)/(4N)$, the condition becomes
$N^2 - 4N + 3 = 0$, giving $N = 1$ (trivial: one family) or
$N = 3$ (physical).
This is a \emph{seventh} $N = 3$ uniqueness criterion,
independent of the Casimir-Koide identity.

A systematic scan of all traceless triples $(a, b, c)$ with
$a + b + c = 0$ and denominators~$\leq 3$ shows that
$(-1, 1/3, 2/3)$ is the \emph{unique} solution giving all
three angles within~$5\%$ of data.
Substituting $N = 2, 3, \ldots, 10$ into
Eqs.~(\ref{eq:t12})--(\ref{eq:t13}) gives $\chi^2 = 0.5$
only at $N = 3$; the next-best is $N = 4$ with $\chi^2 = 218$.

\paragraph{The $1/N$ expansion and the CKM--PMNS hierarchy.}

All mixing parameters are rational functions of~$N$:
\begin{center}
\begin{tabular}{llll}
\toprule
Parameter & Formula & Value & Order \\
\midrule
$\delta$ & $2/N^2$ & $2/9$ & $O(1/N^2)$ \\
$K$ & $j(j{+}1)/N$ & $2/3$ & $O(1/N)$ \\
$|V_{us}|$ & $2/N^2$ & $2/9$ & $O(1/N^2)$ \\
$A$ & $N^2/(N^2{+}2)$ & $9/11$ & $1{-}O(1/N^2)$ \\
$|V_{cb}|$ & $4/(N^4{+}2N^2)$ & $4/99$ & $O(1/N^4)$ \\
$\theta_{12}{-}\pi/4$ & $-2/N^2{+}4/N^4$ & & $O(1/N^2)$ \\
$\theta_{23}{-}\pi/4$ & $2/N^3$ & $2/27$ & $O(1/N^3)$ \\
$\theta_{13}$ & $4/N^3$ & $4/27$ & $O(1/N^3)$ \\
\bottomrule
\end{tabular}
\end{center}

CKM elements scale as $\delta^{|i-j|}$ ($O(1/N^{2|i-j|})$);
PMNS corrections scale as $1/N^{|i-j|-1}$ ($O(1/N)$).
The CKM hierarchy $|V_{us}| \gg |V_{cb}| \gg |V_{ub}|$
follows from the geometric series in~$\delta$;
the near-maximality of the PMNS angles follows from
their being $O(1/N)$ corrections to~$\pi/4$.
Both originate from the same $\SU(3)_F$ but at different
orders: the CKM is a second-order effect ($\delta = O(1/N^2)$)
while the PMNS is first-order ($1/N$).

%======================================================================
\section{The $\mathcal{N}{=}2$ structure}
\label{sec:N2}
%======================================================================

The algebraic structures discovered in the preceding
sections---the Koide formula, the Cabibbo identity, the
bion relations, the democratic resolvent---all point to a
common origin in $\mathcal{N}{=}2$ SQCD.
The Koide formula is the trigonometric identity of the
Cartan torus; the Cabibbo identity is the Cartan angle
itself; the bion relations are the $\mathbb{Q}[\sqrt{3}]$
arithmetic of the seed; and the democratic resolvent is
the propagator on the zero-mode of the Cartan subalgebra.
All of these are properties of the $\mathcal{N}{=}2$
Coulomb branch.

The three hypermultiplets of $\SU(3)_F$ with $\Nf = \Nc = 3$
descend from $\mathcal{N}{=}2$ SQCD.  The UV theory contains,
in addition to the hypermultiplets, an $\mathcal{N}{=}2$ vector
multiplet with gauge bosons, two sets of gauginos
($\lambda^{1,a}$, $\lambda^{2,a}$), and an adjoint complex
scalar~$\Phi^a$ ($a = 1,\ldots,8$).
The Coulomb branch is parameterised by the gauge-invariant
moduli $u_2 = \frac{1}{2}\Tr(\Phi^2)$ and $u_3 = \det(\Phi)$.

The adjoint satisfies the Cayley--Hamilton equation
\begin{equation}
  \Phi^3 - u_2\,\Phi - u_3 = 0\,,
  \label{eq:CH}
\end{equation}
since $\Tr(\Phi) = 0$.  With the identifications
$u_2 = \sum m_i^2$ and $u_3 = 2\prod m_i$ (from
Eq.~(\ref{eq:trig_identities}): $u_2 = \mu^4 \sum z_i^4$,
$u_3 = 2\mu^6 \prod z_i^2$), the Cayley--Hamilton is
algebraically identical to the cubic $Q_2$ algebra
$R^3 = \Tr(D^2)\,R + 2\Lambda^6\,I$ discovered in the
meson operator analysis.
Here $R_{ij} = \langle M_k\rangle$ is the complementary meson
VEV operator and $D = \mathrm{diag}(m_1, m_2, m_3)$.
The cubic is the \emph{residual} $\mathcal{N}{=}2$ structure
after confinement.

For the lepton masses, the eigenvalues of~$\Phi$ are
$\phi_1 = +1780$, $\phi_2 = -0.06$, $\phi_3 = -1780\MeV$,
satisfying $\sum\phi_i = 0$ (traceless) and
$\sum\phi_i^2 = 2\sum m_i^2$ (verified numerically).
The near-vanishing of~$\phi_2$ ($\approx m_e/8$) signals
proximity to an Argyres--Douglas-like point where
$\det(\Phi) \to 0$.

\paragraph{$\mathcal{N}{=}2 \to \mathcal{N}{=}1$ breaking.}

The breaking is effected by an adjoint mass term
$W = \mu_\Phi\,\Tr(\Phi^2)$.  For $\mu_\Phi \to \infty$
the adjoint decouples and the dynamical Cayley--Hamilton
becomes a static algebraic constraint on the meson
eigenvalues.  The effective superpotential becomes
\begin{equation}
  W_{\mathrm{eff}} = \frac{\det M}{\Lambda^3} + m_i M^i_i
  - \frac{1}{2\mu_\Phi}\,\Tr(M^2)\,,
\end{equation}
where the last term is the leading correction from integrating
out~$\Phi$.  The adjoint mass is determined from the CKM
Wolfenstein~$A$ parameter:
$A = 1/(1+\delta)$ gives
$\mu_\Phi = m_\tau(1 + 1/\delta) = 11\,m_\tau/2 \approx
9.8\GeV$ (Sec.~\ref{sec:mixing}).

\paragraph{The Seiberg--Witten curve.}

The SW curve for $\SU(3)$ with $\Nf = 3$ is the genus-2
hyperelliptic surface
$y^2 = P(x)^2 - 4\Lambda^3\,Q(x)$,
where $P(x) = x^3 - u_2 x - u_3$ and
$Q(x) = \prod_i(x+m_i)$.
At the physical Koide point, the curve is \emph{smooth} for
all positive~$\Lambda$.
The Koide constraint reduces the 5-dimensional parameter space
$(u_2, u_3, m_1, m_2, m_3)$ to a one-parameter family
in~$\delta$.

In the democratic limit ($\delta \to 0$, $m_1 = m_2 = m_3 = m$),
$P(x) = (x+m)^2(x-2m)$ and the sextic factorises as
$f(x) = (x+m)^3[(x+m)(x-2m)^2 - 4\Lambda^3]$.
A monopole singularity (double root at $x = 0$) appears at
$\Lambda = m = 314\MeV$---the democratic mass equals the
confinement scale.
This singularity is the Argyres--Douglas point of the
democratic vacuum.

%======================================================================
\section{M-theory and branes}
\label{sec:Mtheory}
%======================================================================

The algebraic results of the preceding sections---the
$\SU(3)_F$ confinement, the $\mathcal{N}{=}2$ Coulomb
branch, the Seiberg--Witten curve---all have natural
realisations in string theory.
The Hanany--Witten~\cite{HananyWitten:1997} brane
construction provides the dictionary.
Three D4-branes stretched between two NS5-branes give the
$\SU(3)_F$ gauge group; three D6-branes provide
$\Nf = 3$ flavours.
The relative separation of the NS5-branes sets
$1/g_F^2$; the D4-D6 open strings are the
preons~$Q_i, \tilde Q_i$.
After s-confinement ($\Nf = \Nc$): no magnetic D4-branes
survive; the IR theory is described entirely by D6--D6
strings (mesons) and wrapped D4-branes (baryons).

Quarks are open strings between the colour-D4$_c$ stack and
the D6 flavour branes.  Leptons are D6--D6 strings.  In
M-theory both lift to M5-branes, and the distinction between
``elementary'' and ``composite'' becomes geometric: which
sheet of the M5 carries the M2-brane boundary.

The Seiberg--Witten curve~(\ref{eq:CH}) \emph{is} the M5
shape.  The Cartan angle~$\delta$ is the M5 position on the
M-theory circle.  The Cabibbo angle equals~$\delta$ because
both are set by the same M5 geometry.

The seed angle $\pi/12$ is the angular quantum connecting
the charged-lepton and neutrino sectors:
$\delta_\nu = \delta + \pi/12 = 1/8$ of the
$\mathbb{Z}_3$ period on the Cartan circle.

In this picture, the distinction between ``elementary'' and
``composite'' dissolves: a quark is an open string between
two brane stacks, while a lepton is an open string between
two sheets of the \emph{same} brane stack (the flavour D6s).
Both are equally fundamental---or equally composite---depending
on which brane coordinates one chooses as ``position'' versus
``field.''
The sBootstrap is the statement that this M-theory democracy
extends to the effective field theory: the supercharge~$Q$
that connects mesons to leptons and diquarks to quarks is the
low-energy remnant of the open-string vertex operator that
changes the endpoint boundary conditions from
D6--D6 (meson) to D4$_c$--D6 (quark).

The diquark sector maps to D4$_c$--D4$_c$ strings in the
colour-antisymmetric channel.  The charge-$4/3$ diquark
$[uc]$ corresponds to a string stretching between D4$_c$
branes with both endpoints carrying up-type quantum numbers;
its problematic charge reflects the fact that the $\SO(32)$
embedding requires the full type~I string spectrum, of which
the low-energy QCD composites are only a subset.

The NS5-brane rotation that effects
$\mathcal{N}{=}2 \to \mathcal{N}{=}1$ (the adjoint mass
$\mu_\Phi$ of Sec.~\ref{sec:N2}) has a simple geometric
interpretation: the relative angle between the two NS5-branes
is~$\mu_\Phi L$, where $L$ is the separation in the
compact direction.
At $\mu_\Phi = 11\,m_\tau/2 \approx 9.8\GeV$, the rotation
angle is small ($\sim 10^{-1}$~rad in appropriate units),
consistent with the N=2 structure being ``almost'' preserved
in the confined phase---which is precisely what the Koide
formula requires.

The Hanany--Witten construction also provides a
brane-theoretic interpretation of the Chadha--Nielsen
obstruction.
A local composite supercharge corresponds to a vertex
operator localised at a point on the D6-brane worldvolume.
The factor-of-18 overshoot arises because the vertex
operator probe is sensitive to the full area of the
D6-brane, which scales as $N_c^3$ in the large-$N_c$ limit.
The holographic evasion (Sec.~\ref{sec:tests}) localises the
supercharge wavefunction near the IR wall $z_{\mathrm{IR}}$,
where the warp factor suppresses the anticommutator by
exactly the factor needed to close the algebra.

%======================================================================
\section{Complete predictions}
\label{sec:predictions}
%======================================================================

\begin{table}[h]
\centering
\caption{Predictions from $N = 3$, $\mu = 17.72\;\sqrt{\MeV}$,
  $m_c = 1270\MeV$.  Mean off: $1.1\%$.}
\label{tab:predictions}
\begin{tabular}{llrl}
\toprule
Parameter & Formula & Off \\
\midrule
$m_e,\;m_\mu,\;m_\tau$ & Koide $+$ $\mu$ $+$ $\delta$ &
  $< 0.01\%$ \\
$|V_{us}|$ & $2/9$ & $1.2\%$ \\
$|V_{cb}|$ & $4/99$ & $4.3\%$ \\
$A$ & $9/11$ & $1.0\%$ \\
$\sin^2\theta_{12}$ & $\sin^2(\pi/4-\delta+\delta^2/N)$ & $1.3\%$\\
$\sin^2\theta_{23}$ & $\sin^2(\pi/4+\delta/N)$ & $0.1\%$ \\
$\sin^2\theta_{13}$ & $\sin^2(4/27)$ & $1.9\%$ \\
$m_d/m_s$ & $4/81$ & $1.2\%$ \\
$m_s$ & $91.6\MeV$ & $1.9\%$ \\
$m_b$ & $4141\MeV$ & $1.0\%$ \\
$m_t$ & $171.1\GeV$ & $0.8\%$ \\
$m_\pi^2 - m_\mu^2$ & $f_\pi^2$ & $2.2\%$ \\
$m_\tau$ & $f_\pi m_b/\Lqcd$ & $0.2\%$ \\
$K(-\pi,D_s,B)$ & $2/3$ & $0.1\%$ \\
\bottomrule
\end{tabular}
\end{table}

Not predicted: CKM and PMNS CP phases (two parameters),
Majorana phases (two more), $m_u$ (isospin splitting,
outside the family-symmetry framework), gauge couplings
$g_1, g_2, g_3$ (part of the elementary $\SU(3)_c \times
\SU(2)_L \times \UU(1)_Y$ sector), Higgs self-coupling~$\lambda$,
$\theta_{\mathrm{QCD}}$ (strong CP), and the neutrino
Koide angle~$\delta_\nu$ (fitted to $\delta + \pi/12$ but
not derived).

The statistical assessment, treating $|V_{us}|$ as a
single input (determining~$\delta$), gives
$\chi^2/\mathrm{dof} = 6.2/5 = 1.25$ ($p = 0.28$) for
the remaining five mixing predictions.  The pulls are
$|V_{cb}|$ ($-2.3\sigma$), $A$ ($-0.6\sigma$),
$\sin^2\theta_{12}$ ($-0.3\sigma$),
$\sin^2\theta_{23}$ ($+0.04\sigma$),
$\sin^2\theta_{13}$ ($-0.6\sigma$).
The fit is acceptable.

With $\delta = 2/9$ used without adjustment, $|V_{us}|$
contributes $-4.1\sigma$, reflecting the $1.2\%$ discrepancy
between $2/9 = 0.2222$ and the PDG value~$0.2250$.
This tension may be resolved if the angular
potential~(\ref{eq:potential}) shifts the minimum to
$\delta = 0.2225$.

\paragraph{Parameter reduction.}

The Standard Model has $19$~free parameters with massless
neutrinos (or $28$ with Dirac neutrino masses and PMNS
phases).  The sBootstrap predicts $15$ of these from
$3$~inputs ($N = 3$, $\mu$, $m_c$), reducing the free count
to~$13$.  The $15/28 = 54\%$ reduction is the main
quantitative achievement of the framework.

%======================================================================
\section{Discussion}
\label{sec:discussion}
%======================================================================

\paragraph{The parameter reduction.}

The sBootstrap predicts $15$ of the $28$ free parameters of
the Standard Model from $3$ inputs ($N$, $\mu$, $m_c$),
a $54\%$ reduction.
No GUT achieves this: $\SU(5)$ predicts $\sim 2$ mass
relations ($m_b = m_\tau$ at $M_{\mathrm{GUT}}$, plus gauge
unification); $\SO(10)$ adds $\sim 2$ more.
String compactifications on Calabi--Yau threefolds have
$O(100)$ moduli and predict zero mass values.
The sBootstrap achieves quantitative predictions because
it operates at the QCD scale, where the dynamics is
controlled by a single confinement parameter~$\Lambda_F$
and the mass hierarchy by a single angle~$\delta$.

\paragraph{Relation to the MSSM.}

The sBootstrap is not the MSSM.
In the MSSM, squarks and sleptons are new particles at the
TeV scale; in the sBootstrap, they are the mesons and diquarks
already observed in QCD.
The MSSM has $\sim 120$ free parameters (soft-breaking masses,
$A$-terms, $\mu$-parameter); the sBootstrap has three
($N$, $\mu$, $m_c$).
The two share the same algebraic structure---chiral superfields
with the same gauge quantum numbers---but differ fundamentally
in the breaking mechanism.
MSSM breaking is soft (gravity- or gauge-mediated, at TeV);
sBootstrap breaking is strong (Volkov--Akulov, at $f_\pi$).

The absence of TeV-scale superpartners at the LHC is
therefore not a problem for the sBootstrap: the superpartners
are at the QCD scale, where they have been measured for
decades as hadrons.

\paragraph{Relation to Masiero--Veneziano.}

Masiero and Veneziano~\cite{MasieroVeneziano:1985} constructed
the first prototype model for one leptonic family from
$\SU(3)$ SQCD with $\Nf = \Nc = 3$.
Their model has the same gauge group and matter content as
ours but differs in the treatment of SUSY breaking (they use
generic soft terms, we use Volkov--Akulov), the electroweak
gauging (they gauge $\SU(2)_L \times \UU(1)_Y$ as a subgroup
of the flavour group), and crucially the absence of the Koide
formula (they have no mass predictions).
We extend their framework by adding the charge-squared rule,
the Cartan angle, and the bion relations.

\paragraph{The electron mass.}

The electron mass vanishes at leading order in the Yukawa
ansatz~(\ref{eq:yukawa_ansatz}): $y^e \sim 0$.
This is a feature, not a bug: the tiny electron mass
($m_e/m_\mu = 1/207$) is naturally explained as a higher-order
effect.
The determinant $\det(m_\ell) \propto m_e\,m_\mu\,m_\tau$
is protected by an approximate $\UU(1)_e$ symmetry (the
electron number); its smallness does not require fine-tuning.
The electron mass could arise from electroweak loops mixing
the $(e, \mu, \tau)$ sector, or from sub-leading Yukawa
couplings $y^e_{cd} \sim g^2/16\pi^2$ at one loop.

\paragraph{The charge-$4/3$ problem.}

The single diquark $[uc]$ with charge~$4/3$ has no Standard
Model fermion partner.
In the $\SO(32)$ embedding~\cite{Rivero:2024}, it is part of
the $\mathbf{442}$ states in the adjoint decomposition
$\mathbf{496} = \mathbf{54} \oplus \mathbf{442}$.
These extra states must be projected out by the UV completion.
In the M-theory picture (Sec.~\ref{sec:Mtheory}), the
charge-$4/3$ diquark is a D4$_c$--D4$_c$ string that can form
a neutral condensate~\cite{v159}, removing it from the
low-energy spectrum.

\paragraph{Why $\sqrt{m}$ is primary.}

The charge-squared rule $m = \mu\,z^2$ makes $\sqrt{m}$ the
primary variable.
This is not merely a parametric convenience.
Karasik~\cite{Karasik:2020} showed that the baryon current
in QCD is most naturally written in terms of the ``square root''
$\xi$ of the chiral field ($\xi^2 = U$), not in terms of~$U$
itself.
The standard skyrmion current $B^\mu$ (built from~$U$) and
Karasik's current $H^\mu$ (built from~$\xi$) agree for smooth
configurations but differ for singular ones (the $N_f = 1$
baryon).
In our framework, the same $\xi$-structure appears in the Koide
parametrisation: $\sqrt{m_n} = \mu\,z_n$ is the family charge
in the ``square root space'' of the meson matrix.

%======================================================================
\section{Experimental tests}
\label{sec:tests}
%======================================================================

\paragraph{JUNO: the critical test}

Our $\sin^2\theta_{12} = 0.300$ is $1.3\%$ below the NuFIT
central value~$0.304$.  JUNO will measure this to
$\pm 0.003$, testing our prediction at $\sim 1.3\sigma$.

\paragraph{Lattice: the binary verdict.}

The Chadha--Nielsen calculation~\cite{ChadhaNielsen:1977,v159}
gives an explicit numerical target.
Using the composite supercharge ansatz $Q_\alpha =
\int d^3z\,K(z)\,\epsilon^{abc}\bar q_a\sigma^\mu_{\alpha\dot\beta}
q_b^{\dot\beta}\partial_\mu\phi_{\pi,c}$ with form factor
$K(z) \sim \Lqcd^{-3/2}$ and the QCD condensate
$\langle\bar q q\rangle = -B_0 f_\pi^2$:
\begin{equation}
  \frac{\{Q,\bar Q\}_{\mathrm{local}}}{2\sigma^0 m_\pi}
  \sim N_c\left(\frac{f_\pi^2}{\Lqcd^2}\right)^2
  \sim N_c^3 \approx 27\,.
  \label{eq:CN_factor}
\end{equation}
With the $O(1)$ factor $m_\pi/\Lqcd \approx 0.64$, the
numerical result is~$\sim 18$.
A nonlocal supercharge must suppress this by exactly
$1/18$ to close the algebra.

Five lattice observables, each requiring $\sim 10^7$ core-hours,
provide a binary verdict.
The PCSC residue should satisfy $Z_S = 1.00 \pm 0.05$;
the auxiliary-field VEV should give
$\langle F_\pi\rangle/(f_\pi^2\Lqcd) = -1.0 \pm 0.2$
from the four-quark condensate;
the restored-phase muon mass should vanish,
$m_\mu(T > T_c) < 20\MeV$;
the APE-smeared anticommutator should close,
$\Delta(t)/(2m_\pi) = 1.0 \pm 0.2$;
and a single holographic IR cutoff~$z_{\mathrm{IR}}$ should
simultaneously reproduce $m_\rho = 775\MeV$ and
$m_{a_1} = 1230\MeV$.
Agreement within $5\%$ on any of these establishes nonlocal
hadronic SUSY; a $3\sigma$ deviation on the anticommutator
falsifies all nonlocal evasions.
A realistic timeline is $36$--$48$ months.

\paragraph{$|V_{cb}|$ exclusive/inclusive}

Our $4/99 = 0.0404$ lies between the exclusive ($0.0395$)
and inclusive ($0.0422$) PDG determinations, $1.1\sigma$ from
the exclusive.  Resolution of this long-standing tension
will sharpen the test.

\begin{acknowledgments}
The author thanks M.~Porter for sustained theoretical
discussions (2011--2024) that identified Seiberg duality and
the mesino interpretation as the mechanism underlying the
sBootstrap, C.~Brannen for the circulant mass matrix
parametrisation, and the late M.~Sheppeard (Kea) for
early contributions connecting Koide to representation
theory.
During the preparation of this work the author used Claude
(Anthropic) for algebraic verification, numerical checks,
and drafting assistance.
The author reviewed and edited all content and takes full
responsibility for the publication.
\end{acknowledgments}

\begin{thebibliography}{99}

\bibitem{Miyazawa:1966}
  H.~Miyazawa,
  Prog.\ Theor.\ Phys.\ \textbf{36}, 1266 (1966).

\bibitem{Rivero:2024}
  A.~Rivero,
  Eur.\ Phys.\ J.\ C \textbf{84}, 1058 (2024),
  arXiv:2407.05397.

\bibitem{ChadhaNielsen:1977}
  S.~Chadha and H.~B.~Nielsen,
  Nucl.\ Phys.\ B \textbf{217}, 125 (1977).

\bibitem{VolkovAkulov:1973}
  D.V.~Volkov and V.P.~Akulov,
  Phys.\ Lett.\ B \textbf{46}, 109 (1973).

\bibitem{Seiberg:1995}
  N.~Seiberg,
  Nucl.\ Phys.\ B \textbf{435}, 129 (1995),
  arXiv:hep-ph/9411149.

\bibitem{MasieroVeneziano:1985}
  A.~Masiero and G.~Veneziano,
  Nucl.\ Phys.\ B \textbf{249}, 593 (1985).

\bibitem{Koide:1983}
  Y.~Koide,
  Phys.\ Lett.\ B \textbf{120}, 161 (1983).

\bibitem{Foot:1994}
  R.~Foot,
  arXiv:hep-ph/9402242 (1994).

\bibitem{Karasik:2020}
  A.~Karasik,
  SciPost Phys.\ \textbf{9}, 008 (2020),
  arXiv:2003.07893.

\bibitem{Unsal:2008}
  M.~\"Unsal,
  Phys.\ Rev.\ D \textbf{78}, 065035 (2008).

\bibitem{Raidal:2004}
  M.~Raidal,
  Phys.\ Rev.\ Lett.\ \textbf{93}, 161801 (2004).

\bibitem{Zenczykowski:2012}
  P.~Zenczykowski,
  Phys.\ Rev.\ D \textbf{86}, 117303 (2012).

\bibitem{GattoSartoriTonin:1968}
  R.~Gatto, G.~Sartori, and M.~Tonin,
  Phys.\ Lett.\ B \textbf{28}, 128 (1968).

\bibitem{HananyWitten:1997}
  A.~Hanany and E.~Witten,
  Nucl.\ Phys.\ B \textbf{492}, 152 (1997),
  arXiv:hep-th/9611230.

\bibitem{Brannen:2006}
  C.A.~Brannen,
  \texttt{hep-ph/0606057} (2006).

\end{thebibliography}

\end{document}
